\chapter{第2章：Eisenstein级数与L函数 习题}
\label{ch:ex-ch2}

\difficultykey

\section{基础题 \easy}

\begin{problem}{\easy\ 计算 $\sigma_3(12)$ 与 $G_8(q)$ 首项}
计算 $\sigma_3(12)$，并写出
\[
G_8(q)=\sum_{n\ge1}\sigma_7(n)q^n
\]
的 $q^0,q^1,q^2$ 项。
\end{problem}
\begin{solution}
$12$ 的正因子为 $1,2,3,4,6,12$，故
\[
\sigma_3(12)=1^3+2^3+3^3+4^3+6^3+12^3=1+8+27+64+216+1728=2044.
\]
又
\[
\sigma_7(1)=1,\qquad \sigma_7(2)=1+2^7=129.
\]
因此
\[
G_8(q)=0+q+129q^2+\cdots.
\]
若只从 $n\ge1$ 开始求和，则 $q^0$ 项就是 $0$。
\end{solution}

\begin{problem}{\easy\ Bernoulli数的前几项}
由生成函数
\[
\frac{t}{e^t-1}=\sum_{n\ge0}B_n\frac{t^n}{n!}
\]
求 $B_0,B_2,B_4,B_6$。
\end{problem}
\begin{solution}
展开得
\[
\frac{t}{e^t-1}=1-\frac t2+\frac{t^2}{12}-\frac{t^4}{720}+\frac{t^6}{30240}+\cdots.
\]
与
\[
\sum_{n\ge0}B_n\frac{t^n}{n!}
\]
逐项比较可得
\[
B_0=1,\qquad B_2=\frac16,\qquad B_4=-\frac1{30},\qquad B_6=\frac1{42}.
\]
\end{solution}

\begin{problem}{\easy\ 证明 $E_4^2=E_8$ 并推出卷积恒等式}
证明
\[
E_4^2=E_8,
\]
并由此推出 $\sigma_7$ 与 $\sigma_3$ 的卷积公式。
\end{problem}
\begin{solution}
标准归一化下
\[
E_4=1+240\sum_{n\ge1}\sigma_3(n)q^n,\qquad
E_8=1+480\sum_{n\ge1}\sigma_7(n)q^n.
\]
由于 $\dim M_8=1$，权 $8$ 模形式空间由 $E_8$ 张成；而 $E_4^2$ 也是权 $8$、常数项为 $1$ 的模形式，所以
\[
E_4^2=E_8.
\]
比较 $q^n$ 系数，得
\[
480\sigma_7(n)=480\sigma_3(n)+240^2\sum_{m=1}^{n-1}\sigma_3(m)\sigma_3(n-m).
\]
两边除以 $480$，得到
\[
\sigma_7(n)=\sigma_3(n)+120\sum_{m=1}^{n-1}\sigma_3(m)\sigma_3(n-m).
\]
\end{solution}

\begin{problem}{\easy\ 证明 $G_2$ 的准模性}
证明 $G_2$ 不是模形式而是准模形式，并写出其变换公式。
\end{problem}
\begin{solution}
把
\[
E_2(\tau)=1-24\sum_{n\ge1}\sigma_1(n)q^n
\]
看作 $G_2$ 的归一化形式。它满足
\[
E_2(\tau+1)=E_2(\tau),
\]
但在 $\tau\mapsto -1/\tau$ 下有
\[
E_2\!\left(-\frac1\tau\right)=\tau^2E_2(\tau)+\frac{6}{\pi i}\tau.
\]
更一般地，对
\[
\gamma=\begin{pmatrix}a&b\\ c&d\end{pmatrix}\in \mathrm{SL}_2(\mathbf Z),
\]
有
\[
E_2(\gamma\tau)=(c\tau+d)^2E_2(\tau)+\frac{6}{\pi i}c(c\tau+d).
\]
多出来的线性项说明它不是真正的权 $2$ 模形式，而是权 $2$ 的准模形式。等价地，若用未归一化的 $G_2$，则公式写成
\[
G_2(\gamma\tau)=(c\tau+d)^2G_2(\tau)-2\pi i\,c(c\tau+d).
\]
特别地，取 $\gamma=S=\begin{pmatrix}0&-1\\1&0\end{pmatrix}$，便得到
\[
G_2(-1/\tau)=\tau^2G_2(\tau)-2\pi i\tau.
\]
\end{solution}

\begin{problem}{\easy\ 写出 $E_{12}$ 的 $L$-级数前三项}
对 $E_{12}$ 写出 $L(f,s)$ 的前三项，并说明它与 zeta 函数的关系。
\end{problem}
\begin{solution}
标准归一化下
\[
E_{12}=1+\frac{65520}{691}\sum_{n\ge1}\sigma_{11}(n)q^n.
\]
于是
\[
\sigma_{11}(1)=1,\qquad
\sigma_{11}(2)=1+2^{11}=2049,\qquad
\sigma_{11}(3)=1+3^{11}=177148.
\]
因此
\[
L(E_{12},s)=\frac{65520}{691}\left(1+\frac{2049}{2^s}+\frac{177148}{3^s}+\cdots\right).
\]
又因为
\[
\sum_{n\ge1}\frac{\sigma_{11}(n)}{n^s}=\zeta(s)\zeta(s-11),
\]
所以
\[
L(E_{12},s)=\frac{65520}{691}\,\zeta(s)\zeta(s-11).
\]
\end{solution}

\begin{problem}{\easy\ 证明 $\zeta(2)=\pi^2/6$}
证明
\[
\zeta(2)=\frac{\pi^2}{6}.
\]
\end{problem}
\begin{solution}
Euler 的乘积公式
\[
\frac{\sin \pi x}{\pi x}=\prod_{n\ge1}\left(1-\frac{x^2}{n^2}\right)
\]
两边展开并比较 $x^2$ 项。左边为
\[
1-\frac{\pi^2x^2}{6}+\cdots,
\]
右边为
\[
1-\left(\sum_{n\ge1}\frac1{n^2}\right)x^2+\cdots.
\]
故
\[
\sum_{n\ge1}\frac1{n^2}=\frac{\pi^2}{6},
\]
即
\[
\zeta(2)=\frac{\pi^2}{6}.
\]
\end{solution}

\begin{problem}{\easy\ 计算若干权的 $\dim S_k$}
计算
\[
\dim S_k\qquad (k=12,14,16,18,20,22).
\]
\end{problem}
\begin{solution}
在 $\mathrm{SL}_2(\mathbf Z)$ 上，
\[
S_k=\Delta\cdot M_{k-12}\qquad (k\ge12,\ k\text{ 偶}).
\]
因此
\[
\dim S_k=\dim M_{k-12}.
\]
而
\[
\dim M_0=1,\ \dim M_2=0,\ \dim M_4=1,\ \dim M_6=1,\ \dim M_8=1,\ \dim M_{10}=1.
\]
所以
\[
\dim S_{12}=1,\quad
\dim S_{14}=0,\quad
\dim S_{16}=1,\quad
\dim S_{18}=1,\quad
\dim S_{20}=1,\quad
\dim S_{22}=1.
\]
\end{solution}

\begin{problem}{\easy\ 证明 $\sum \sigma_k(n)n^{-s}=\zeta(s)\zeta(s-k)$}
证明在收敛区域内有
\[
\sum_{n\ge1}\frac{\sigma_k(n)}{n^s}=\zeta(s)\zeta(s-k).
\]
\end{problem}
\begin{solution}
由
\[
\sigma_k(n)=\sum_{d\mid n}d^k
\]
可得
\[
\sum_{n\ge1}\frac{\sigma_k(n)}{n^s}
=
\sum_{n\ge1}\sum_{d\mid n}\frac{d^k}{n^s}.
\]
令 $n=dm$，则
\[
\sum_{n\ge1}\frac{\sigma_k(n)}{n^s}
=
\sum_{d,m\ge1}\frac{d^k}{(dm)^s}
=
\left(\sum_{d\ge1}d^{k-s}\right)\left(\sum_{m\ge1}m^{-s}\right).
\]
故
\[
\sum_{n\ge1}\frac{\sigma_k(n)}{n^s}
=
\zeta(s-k)\zeta(s).
\]
当 $\Re(s)>k+1$ 时上式绝对收敛。
\end{solution}

\begin{problem}{\easy\ 说明为什么 $\Delta=q\prod_{n\ge1}(1-q^n)^{24}$}
说明
\[
\Delta(\tau)=q\prod_{n\ge1}(1-q^n)^{24}
\]
为什么是权 $12$ 的尖点形式。
\end{problem}
\begin{solution}
Dedekind eta 函数为
\[
\eta(\tau)=q^{1/24}\prod_{n\ge1}(1-q^n),
\]
它是权 $1/2$ 的模形式（带乘子）。因此
\[
\eta(\tau)^{24}=q\prod_{n\ge1}(1-q^n)^{24}
\]
是权 $12$ 的模形式，而且乘子已消失。
又它的 $q$-展开首项为 $q$，故在无穷远尖点处消失，是 cusp form。由于 $S_{12}$ 一维，这个首项为 $q$ 的权 $12$ 尖点形式只能是 $\Delta$，所以
\[
\Delta(\tau)=\eta(\tau)^{24}=q\prod_{n\ge1}(1-q^n)^{24}.
\]
\end{solution}

\begin{problem}{\easy\ 验证 $E_4$ 的函数方程}
验证 $E_4$ 的完成 $L$-函数满足关于 $s\leftrightarrow 4-s$ 的函数方程。
\end{problem}
\begin{solution}
考虑去掉常数项后的 Mellin 变换
\[
\Lambda^\ast(E_4,s):=\int_0^\infty \bigl(E_4(iy)-1-y^{-4}\bigr)y^{s-1}\,dy.
\]
它收敛，并代表 $E_4$ 的正则化完成 $L$-函数。由模性
\[
E_4\!\left(-\frac1\tau\right)=\tau^4E_4(\tau),
\]
取 $\tau=iy$ 得
\[
E_4(i/y)=y^4E_4(iy).
\]
把积分在 $(0,1)$ 上作代换 $y\mapsto 1/y$，便得到
\[
\Lambda^\ast(E_4,s)=\Lambda^\ast(E_4,4-s).
\]
这就是 $E_4$ 的函数方程。
\end{solution}

\begin{problem}{\easy\ 原始 Dirichlet 特征的 Euler 乘积与延拓}
设 $\chi$ 是原始 Dirichlet 特征。写出 $L(\chi,s)$ 的 Euler 乘积，并说明它如何解析延拓到 $\Re(s)>0$。
\end{problem}
\begin{solution}
在 $\Re(s)>1$ 时，
\[
L(\chi,s)=\sum_{n\ge1}\frac{\chi(n)}{n^s}
=\prod_p\left(1-\chi(p)p^{-s}\right)^{-1}.
\]
若 $\chi$ 非平凡且原始，设 $a=0$ 或 $1$ 分别表示 $\chi(-1)=1$ 或 $-1$，则完成函数
\[
\Lambda(\chi,s)=\left(\frac{q}{\pi}\right)^{(s+a)/2}\Gamma\!\left(\frac{s+a}{2}\right)L(\chi,s)
\]
满足函数方程
\[
\Lambda(\chi,s)=\varepsilon(\chi)\Lambda(\overline{\chi},1-s),
\qquad |\varepsilon(\chi)|=1.
\]
因此 $L(\chi,s)$ 可延拓为整函数；当然也就延拓到了 $\Re(s)>0$。唯一的例外是平凡特征，它给出 $\zeta(s)$，在 $s=1$ 有简单极点。
\end{solution}

\begin{problem}{\easy\ 负奇整数处的 zeta 值}
用 Bernoulli 数表示 $\zeta(-1),\zeta(-3),\zeta(-5)$，并解释为什么 $\zeta(-2m)=0$ 称为平凡零点。
\end{problem}
\begin{solution}
公式
\[
\zeta(1-n)=-\frac{B_n}{n}\qquad (n\ge1)
\]
给出
\[
\zeta(-1)=-\frac{B_2}{2}=-\frac1{12},
\qquad
\zeta(-3)=-\frac{B_4}{4}=\frac1{120},
\qquad
\zeta(-5)=-\frac{B_6}{6}=-\frac1{252}.
\]
又因
\[
B_{2m+1}=0\qquad (m\ge1),
\]
所以
\[
\zeta(-2m)=\zeta\bigl(1-(2m+1)\bigr)=-\frac{B_{2m+1}}{2m+1}=0.
\]
这些零点直接来自 Bernoulli 数的消失，位置又很显然，所以称为平凡零点。
\end{solution}

\section{进阶题 \medium}

\begin{problem}{\medium\ 证明 $L(E_k,s)=\zeta(s)\zeta(s-k+1)$}
证明 Eisenstein 级数的 Dirichlet 级数满足
\[
L(E_k,s)=\zeta(s)\zeta(s-k+1).
\]
\end{problem}
\begin{solution}
若把 Eisenstein 级数的非零次 Fourier 系数记为
\[
a_n=\sigma_{k-1}(n),
\]
则
\[
L(E_k,s)=\sum_{n\ge1}\frac{a_n}{n^s}
=\sum_{n\ge1}\frac{\sigma_{k-1}(n)}{n^s}.
\]
由上一节的恒等式，立刻得到
\[
L(E_k,s)=\zeta(s)\zeta\bigl(s-(k-1)\bigr)=\zeta(s)\zeta(s-k+1).
\]
\textbf{注。} 若采用标准归一化
\[
E_k=1-\frac{2k}{B_k}\sum_{n\ge1}\sigma_{k-1}(n)q^n,
\]
则右边还要乘上同一个常数 $-2k/B_k$。
\end{solution}

\begin{problem}{\medium\ 证明尖点形式 $L(f,s)$ 的绝对收敛}
设 $f(\tau)=\sum_{n\ge1}a_nq^n\in S_k$。证明 $L(f,s)=\sum a_nn^{-s}$ 在 $\Re(s)>k$ 时绝对收敛。
\end{problem}
\begin{solution}
由 Fourier 系数公式
\[
a_n=e^{2\pi ny}\int_0^1 f(x+iy)e^{-2\pi inx}\,dx.
\]
又由模性可知，当 $y\to0^+$ 时
\[
f(x+iy)=O(y^{-k/2}),
\]
取 $y=1/n$，便得到
\[
a_n=O(n^{k/2}).
\]
于是当 $\Re(s)=\sigma>k$ 时，
\[
\sum_{n\ge1}|a_n|n^{-\sigma}
\ll
\sum_{n\ge1}n^{k/2-\sigma}.
\]
因为 $\sigma>k\ge k/2+1$，故右端收敛，所以 $L(f,s)$ 在 $\Re(s)>k$ 绝对收敛。
\end{solution}

\begin{problem}{\medium\ Mellin变换导出的 Hecke 函数方程}
用 Mellin 变换说明 Hecke 型函数方程的来源。
\end{problem}
\begin{solution}
设 $f$ 是权 $k$、level $N$ 的 cusp form，并满足 Fricke 关系
\[
f\big|_k W_N=\varepsilon\,\overline{f},\qquad |\varepsilon|=1.
\]
定义完成 $L$-函数
\[
\Lambda(f,s)=N^{s/2}(2\pi)^{-s}\Gamma(s)L(f,s)
=\int_0^\infty f\!\left(\frac{iy}{\sqrt N}\right)y^{s-1}\,dy.
\]
在积分中作代换 $y\mapsto 1/y$，再用 Fricke 关系，可得
\[
\Lambda(f,s)=\varepsilon\, i^k\,\Lambda(\overline f,k-s).
\]
若 $f$ 系数实，则化为
\[
\Lambda(f,s)=\varepsilon\, i^k\,\Lambda(f,k-s).
\]
这就是 Hecke 函数方程。
\end{solution}

\begin{problem}{\medium\ Poincar\'e级数与 Petersson 内积}
说明 Poincar\'e 级数 $P_m$ 属于 $S_k$，并计算它与任意 $f\in S_k$ 的 Petersson 内积。
\end{problem}
\begin{solution}
定义
\[
P_m(\tau)=\sum_{\gamma\in \Gamma_\infty\backslash \Gamma}
\frac{e^{2\pi i m\gamma\tau}}{j(\gamma,\tau)^k}
\qquad (m\ge1,\ k>2).
\]
它按构造满足权 $k$ 的模变换；又因为 $m\ge1$，在尖点处有指数衰减，所以
\[
P_m\in S_k.
\]
若
\[
f(\tau)=\sum_{n\ge1}a_nq^n\in S_k,
\]
则 unfolding 计算给出
\[
\langle f,P_m\rangle
=
\frac{\Gamma(k-1)}{(4\pi m)^{k-1}}\,a_m.
\]
因此 $P_m$ 的作用是“抽取第 $m$ 个 Fourier 系数”。
\end{solution}

\begin{problem}{\medium\ 说明为什么 $L(E_4,1)=\pi^4/90$}
从 zeta 因子分解解释为什么会出现
\[
\frac{\pi^4}{90}.
\]
\end{problem}
\begin{solution}
由
\[
L(E_4,s)=\zeta(s)\zeta(s-3)
\]
可见 Eisenstein 情形在 $s=1$ 会遇到 $\zeta(s)$ 的极点，所以通常要先把这一个极点因子剥去。剥去之后剩下的有限部分就是
\[
\zeta(4)=\frac{\pi^4}{90}.
\]
因此题目中的
\[
L(E_4,1)=\frac{\pi^4}{90}
\]
应理解为 Eisenstein 极点去掉后的正规化 special value。

\textbf{注。} 真正未正规化的 $L(E_4,s)$ 在 $s=1$ 并不直接取有限值。
\end{solution}

\begin{problem}{\medium\ 椭圆曲线好素数处的 Euler 因子}
设 $E/\mathbf Q$ 是椭圆曲线，$p\nmid N$ 为好素数。证明局部 Euler 因子为
\[
(1-a_pp^{-s}+p^{1-2s})^{-1},
\qquad
a_p=p+1-\#E(\mathbf F_p).
\]
\end{problem}
\begin{solution}
对好素数 $p$，局部 zeta 函数满足 Weil 形式
\[
Z(E/\mathbf F_p,T)=\frac{1-a_pT+pT^2}{(1-T)(1-pT)}.
\]
其中分子是 Frobenius 在一维上同调上的特征多项式。设其根为 $\alpha_p,\beta_p$，则
\[
\alpha_p+\beta_p=a_p,\qquad \alpha_p\beta_p=p.
\]
于是局部 $L$-因子取分子的倒数：
\[
L_p(E,s)=\frac1{(1-\alpha_pp^{-s})(1-\beta_pp^{-s})}
=\frac1{1-a_pp^{-s}+p^{1-2s}}.
\]
即得所求。
\end{solution}

\begin{problem}{\medium\ Hecke本征值的代数整数性}
证明归一化 Hecke 本征形式的本征值 $a_n$ 都是代数整数。
\end{problem}
\begin{solution}
Hecke 算子 $T_n$ 保持模形式空间中的整数格，例如保持整系数 $q$-展开所张成的有限生成模。因此 $T_n$ 在某组基下可以写成整系数矩阵，所以它的特征多项式属于 $\mathbf Z[x]$，且首一。
若 $f$ 是归一化 Hecke 本征形式，则
\[
T_nf=a_nf.
\]
因此 $a_n$ 是整系数首一多项式的根，从而是代数整数。
\end{solution}

\begin{problem}{\medium\ Eisenstein空间与 cusp空间的正交性}
说明为什么 Eisenstein 子空间与 cusp 子空间在 Petersson 内积下正交。
\end{problem}
\begin{solution}
只需验证
\[
\langle E_k,f\rangle=0\qquad (f\in S_k).
\]
把 $E_k$ 写成对 $\Gamma_\infty\backslash\Gamma$ 的和后做 unfolding，得到
\[
\langle E_k,f\rangle
=
\int_0^\infty\int_0^1 \overline{f(x+iy)}\,y^{k-2}\,dx\,dy.
\]
但 cusp form 的常数项为 $0$，所以对每个 $y>0$ 都有
\[
\int_0^1 f(x+iy)\,dx=0.
\]
故上式为 $0$。因此 Eisenstein 子空间与 cusp 子空间正交。
\end{solution}

\begin{problem}{\medium\ 用 $\zeta(12)$ 表示 $\langle \Delta,\Delta\rangle$}
给出把 Petersson 范数 $\langle \Delta,\Delta\rangle$ 表成 $\zeta(12)$ 的一个标准表达式。
\end{problem}
\begin{solution}
Rankin--Selberg 方法给出
\[
\langle \Delta,\Delta\rangle
=
\frac{(4\pi)^{11}}{11!}\cdot\frac{\pi}{3}\,
L(1,\operatorname{Sym}^2\Delta).
\]
又
\[
\zeta(12)=\frac{691}{638512875}\pi^{12},
\qquad
\pi^{12}=\frac{638512875}{691}\zeta(12).
\]
因此可把上式改写为
\[
\langle \Delta,\Delta\rangle
=
\frac{4^{11}\cdot 638512875}{3\cdot 11!\cdot 691}\,
\zeta(12)\,L(1,\operatorname{Sym}^2\Delta).
\]
这就把 $\langle \Delta,\Delta\rangle$ 明确写成了“$\zeta(12)$ 乘以一个对称平方特殊值”的形式。
\end{solution}

\begin{problem}{\medium\ 说明导体 $11$ 曲线的模性并验证到 $p=7$}
设
\[
E_{11}: y^2+y=x^3-x^2-10x-20.
\]
说明它的模性，并验证 $p=2,3,5,7$ 时的 Fourier 系数。
\end{problem}
\begin{solution}
模性定理说明：每条定义在 $\mathbf Q$ 上的椭圆曲线都对应一个权 $2$ 新形式 $f_E=\sum a_nq^n$，使得
\[
L(E,s)=L(f_E,s).
\]
导体 $11$ 的这条曲线对应的新形式是
\[
f_{11}=q-2q^2-q^3+2q^4+q^5+2q^6-2q^7+\cdots.
\]
在好素数处
\[
a_p=p+1-\#E(\mathbf F_p).
\]
直接计数得
\[
\#E(\mathbf F_2)=5,\quad
\#E(\mathbf F_3)=5,\quad
\#E(\mathbf F_5)=5,\quad
\#E(\mathbf F_7)=10.
\]
故
\[
a_2=2+1-5=-2,\quad
a_3=3+1-5=-1,\quad
a_5=5+1-5=1,\quad
a_7=7+1-10=-2.
\]
这与 $f_{11}$ 的系数完全一致，因此在 $p\le7$ 的范围内验证了模性对应。
\end{solution}

\begin{problem}{\medium\ zeta在 $\Re(s)>0$ 的解析延拓公式}
写出一个把 $\zeta(s)$ 延拓到 $\Re(s)>0$ 的公式，并说明原因。
\end{problem}
\begin{solution}
交错级数
\[
\eta(s)=\sum_{n\ge1}\frac{(-1)^{n-1}}{n^s}
\]
在 $\Re(s)>0$ 收敛，并定义全纯函数。又在 $\Re(s)>1$ 时，
\[
\eta(s)=(1-2^{1-s})\zeta(s).
\]
因此
\[
\zeta(s)=\frac{\eta(s)}{1-2^{1-s}}
=\frac{1}{1-2^{1-s}}\sum_{n\ge1}\frac{(-1)^{n-1}}{n^s}.
\]
这就把 $\zeta(s)$ 延拓到了 $\Re(s)>0$；唯一真正的奇点是 $s=1$ 的简单极点。
\end{solution}

\begin{problem}{\medium\ 用 theta 与 Poisson 求和证明 $\xi(s)=\xi(1-s)$}
利用 theta 函数和 Poisson 求和说明 Riemann 函数方程。
\end{problem}
\begin{solution}
设
\[
\theta(t)=\sum_{n\in \mathbf Z}e^{-\pi n^2t}.
\]
Poisson 求和给出
\[
\theta(t)=t^{-1/2}\theta(1/t).
\]
又在 $\Re(s)>1$ 时，
\[
\pi^{-s/2}\Gamma\!\left(\frac s2\right)\zeta(s)
=
\frac12\int_0^\infty (\theta(t)-1)t^{s/2-1}\,dt.
\]
记
\[
\xi(s):=\pi^{-s/2}\Gamma\!\left(\frac s2\right)\zeta(s).
\]
把积分拆成 $(0,1)$ 与 $(1,\infty)$，并在 $(0,1)$ 上用 $t\mapsto 1/t$ 及 theta 变换公式，可得
\[
\xi(s)=\frac12\int_1^\infty (\theta(t)-1)\bigl(t^{s/2-1}+t^{(1-s)/2-1}\bigr)\,dt.
\]
右端在 $s\leftrightarrow 1-s$ 下不变，因此
\[
\xi(s)=\xi(1-s).
\]
\end{solution}

\begin{problem}{\medium\ 陈述广义Ramanujan与GRH及其和椭圆曲线的关系}
陈述广义 Ramanujan 猜想与广义 Riemann 猜想，并说明它们与椭圆曲线 $L$-函数的关系。
\end{problem}
\begin{solution}
广义 Ramanujan 猜想说：未分歧处的归一化 Satake 参数都应落在单位圆上。对权 $2$ 模形式，这等价于
\[
|a_p|\le 2\sqrt p.
\]
对椭圆曲线，这正是 Hasse 界。

广义 Riemann 猜想说：适当完成后的自守 $L$-函数，其所有非平凡零点都落在临界线上。对椭圆曲线的 Hasse--Weil $L(E,s)$，这意味着非平凡零点应位于
\[
\Re(s)=1
\]
（若改用居中变量，则是 $\Re(s)=1/2$）。

二者分别控制局部 Euler 因子的大小和全局零点的分布，因此都直接影响椭圆曲线点数波动、素数和估计以及 BSD 型问题。
\end{solution}

\section{挑战题 \hard}

\begin{problem}{\hard\ 证明偶权 $k\ge12$ 时 $\dim S_k=\dim M_k-1$}
对偶权 $k\ge12$，证明
\[
\dim S_k=\dim M_k-1.
\]
\end{problem}
\begin{solution}
在 $\mathrm{SL}_2(\mathbf Z)$ 上，偶权 $k\ge4$ 的 Eisenstein 子空间是一维，由 $E_k$ 张成；而 cusp 子空间 $S_k$ 正好是它的补空间，所以
\[
M_k=\mathbf C E_k\oplus S_k.
\]
因此
\[
\dim M_k=1+\dim S_k,
\qquad
\dim S_k=\dim M_k-1.
\]
对所有偶权 $k\ge12$ 均成立。
\end{solution}

\begin{problem}{\hard\ 陈述 Petersson 迹公式并化简 $m=n=1$}
陈述 Petersson 迹公式，并把它化简到 $m=n=1$。
\end{problem}
\begin{solution}
设 $\mathcal B_k$ 是 $S_k(\mathrm{SL}_2(\mathbf Z))$ 的一组正交归一 Hecke 基，则
\[
\sum_{f\in\mathcal B_k}
\frac{\Gamma(k-1)}{(4\pi\sqrt{mn})^{k-1}}
a_f(m)\overline{a_f(n)}
=
\delta_{m,n}
+
2\pi i^{-k}\sum_{c\ge1}\frac{S(m,n;c)}{c}
J_{k-1}\!\left(\frac{4\pi\sqrt{mn}}{c}\right),
\]
其中 $S(m,n;c)$ 是 Kloosterman 和，$J_{k-1}$ 是 Bessel 函数。
当 $m=n=1$ 时，
\[
\sum_{f\in\mathcal B_k}
\frac{\Gamma(k-1)}{(4\pi)^{k-1}}|a_f(1)|^2
=
1+
2\pi i^{-k}\sum_{c\ge1}\frac{S(1,1;c)}{c}
J_{k-1}\!\left(\frac{4\pi}{c}\right).
\]
若取 Hecke 归一化 $a_f(1)=1$，左边就是一组显式权重下的基元素个数之和。
\end{solution}

\begin{problem}{\hard\ 用解析类数公式计算 $h(-23)=3$}
用解析类数公式证明
\[
h(-23)=3.
\]
\end{problem}
\begin{solution}
对虚二次域 $\mathbf Q(\sqrt{-23})$，判别式为 $-23$，单位根个数为 $2$。解析类数公式给出
\[
L(1,\chi_{-23})=\frac{2\pi}{2\sqrt{23}}\,h(-23)=\frac{\pi}{\sqrt{23}}h(-23).
\]
另一方面，对这个二次特征可用有限和公式计算
\[
L(1,\chi_{-23})=\frac{3\pi}{\sqrt{23}}.
\]
于是比较两式得到
\[
\frac{\pi}{\sqrt{23}}h(-23)=\frac{3\pi}{\sqrt{23}},
\]
所以
\[
h(-23)=3.
\]
\end{solution}

\begin{problem}{\hard\ 陈述强BSD并说明 $L(1,E)\neq0$ 与秩 $0$ 的关系}
陈述强 Birch--Swinnerton-Dyer 猜想，并解释为什么 $L(1,E)\neq0$ 预示秩为 $0$；同时说明已知结果。
\end{problem}
\begin{solution}
强 BSD 猜想断言
\[
\operatorname{ord}_{s=1}L(E,s)=\operatorname{rank}E(\mathbf Q),
\]
且若 $r=\operatorname{rank}E(\mathbf Q)$，则
\[
\frac{L^{(r)}(E,1)}{r!}
=
\frac{\Omega_E\,\operatorname{Reg}_E\,|\Sha(E)|\,\prod_pc_p}
{|E(\mathbf Q)_{\mathrm{tors}}|^2}.
\]
因此若
\[
L(1,E)\neq0,
\]
则解析秩为 $0$；按 BSD 便应推出代数秩也为 $0$。这在猜想层面给出“非零 special value $\Rightarrow$ rank $0$”。

对 $\mathbf Q$ 上的模椭圆曲线，Gross--Zagier 与 Kolyvagin 的结果进一步表明：若 $L(1,E)\neq0$，则的确有
\[
\operatorname{rank}E(\mathbf Q)=0,
\]
并且 $\Sha(E)$ 有限。
\end{solution}

\begin{problem}{\hard\ 讨论 $E_{11}$ 的部分乘积与 BSD 证据}
讨论
\[
\prod_{p\le100}\frac{\#E(\mathbf F_p)}{p}
\]
对导体 $11$ 的椭圆曲线 $E_{11}$ 给出的 BSD 数值证据。
\end{problem}
\begin{solution}
对好素数 $p$，
\[
\frac{\#E(\mathbf F_p)}{p}
=
1-\frac{a_p}{p}+\frac1p
=
L_p(E,1)^{-1}.
\]
因此这个部分乘积可看作 $L(E,1)^{-1}$ 的一个截断近似。对
\[
E_{11}: y^2+y=x^3-x^2-10x-20
\]
直接计算到 $p\le100$，得到
\[
\prod_{p\le100}\frac{\#E(\mathbf F_p)}{p}\approx 6.4878.
\]
于是其倒数约为
\[
0.1541.
\]
这说明截断 Euler 乘积并没有向 $0$ 衰减，而是表现为趋向某个正的有限值；这与
\[
L(E,1)\neq0
\]
的预期一致，因此支持解析秩为 $0$，也就支持 BSD 所预言的
\[
\operatorname{rank}E_{11}(\mathbf Q)=0.
\]

\textbf{注。} 这只是数值证据：收敛很慢，坏素数的局部因子也需单独处理，但定性上它确实支持 rank $0$ 的 BSD 图景。
\end{solution}
